\begin{center}
\begin{tikzpicture}[x=.88cm,y=.88cm,font=\small]
  \fill[paper] (0,0) rectangle (11.4,5.8);

  \fill[dynasty!19] (.75,3.55) -- (3.2,3.45) -- (3.45,5.15) -- (.95,5.3) -- cycle;
  \node[align=center] at (2.1,4.4) {晋\\北方强国};
  \fill[source!15] (.95,1.15) -- (3.15,1.1) -- (3.3,2.55) -- (1.1,2.75) -- cycle;
  \node[align=center] at (2.15,1.9) {晋的从属网络\\随会盟与征伐变动};
  \fill[timeline!20] (8.15,3.55) -- (10.6,3.45) -- (10.8,5.15) -- (8.4,5.3) -- cycle;
  \node[align=center] at (9.55,4.4) {楚\\南方强国};
  \fill[source!15] (8.3,1.15) -- (10.5,1.1) -- (10.65,2.55) -- (8.45,2.75) -- cycle;
  \node[align=center] at (9.55,1.9) {楚的从属网络\\随会盟与征伐变动};
  \fill[timeline!22] (4.65,2.25) -- (6.75,2.15) -- (6.95,3.8) -- (4.85,3.95) -- cycle;
  \node[align=center] at (5.8,3.15) {宋\\向戌居间};
  \fill[dynasty!13] (4.0,.55) -- (7.6,.5) -- (7.75,1.45) -- (4.15,1.55) -- cycle;
  \node[align=center] at (5.85,1.02) {前546年弭兵之会\\暂时并列两套关系};

  \draw[-{Stealth[length=2mm]},ultra thick,color=dynasty]
    (2.7,4.0) .. controls (4.1,4.35) and (4.7,3.85) .. (5.0,3.55);
  \draw[-{Stealth[length=2mm]},ultra thick,color=timeline]
    (8.9,4.0) .. controls (7.5,4.35) and (6.9,3.85) .. (6.6,3.55);
  \draw[{Stealth[length=2mm]}-{Stealth[length=2mm]},thick,dashed,color=source]
    (3.2,2.1) .. controls (4.2,2.45) and (4.55,2.8) .. (4.8,2.95);
  \draw[{Stealth[length=2mm]}-{Stealth[length=2mm]},thick,dashed,color=source]
    (8.35,2.1) .. controls (7.4,2.45) and (7.05,2.8) .. (6.8,2.95);
  \draw[-{Stealth[length=2mm]},thick,color=source]
    (5.8,2.2) -- (5.85,1.55);

  \node[text=dynasty,above] at (3.85,4.2) {征发、竞争};
  \node[text=timeline,above] at (7.75,4.2) {征发、竞争};
  \node[text=source,align=center] at (3.65,2.85) {使者、盟约\\往来方向};
  \node[text=source,align=center] at (7.95,2.85) {使者、盟约\\往来方向};
  \node[text=source,align=center,font=\footnotesize] at (5.8,5.2)
    {弭兵不是统一指挥体系，也没有常设执行机构};
\end{tikzpicture}

\smallskip
\small\textit{图\quad 宋地弭兵与晋楚双重网络示意：以宋向戌的居间往来说明，前546年的安排暂时承认晋、楚各自的从属网络，而非建立统一政府。参考《春秋·襄公二十七年》《左传·襄公二十七年》及《国语·宋语》；谈判言辞、从属国名单和约束范围主要由后出叙事展开，仍须辨析。图中关系、地点和使者路线均为约略示意，不表示精确疆界、实际行程或稳定的隶属秩序。}
\end{center}
